Пусть $L \supset F$ -- конечное расширение полей. Тогда любой гомоморфизм
$\varphi: L \to L$, сохраняющий $F$ (то есть для любого $a \in F$ выполнено
равенство $\varphi(a) = a$) является изоморфизмом на себя, т.е. автоморфизмом.

Докажем, что $\varphi$ инъективен.

Предположим, что $Ker \varphi \neq \{0\}$. 
Тогда $\exists a \in Ker \varphi\ :\ a \neq 0$.
Но тогда $\varphi(1) = \varphi(a \cdot a^{-1}) = 
\varphi(a) \cdot \varphi(a^{-1}) = 0$. Противоречие с тем, что
$\varphi$ -- гомоморфизм.

Теперь докажем, что $\varphi$ сюръективен.

Так как $L$ -- линейное пространство над $F$, то можно считать, что
$L \cong F^k$, где $k = [L : K]$. 
Тогда $\varphi: F^k \to F^k$ -- линейный оператор.
Из равенства
$dim\ Ker \varphi + dim\ Im \varphi = k$ получаем, что
$Im \varphi = F^k$, то есть $\varphi$ сюръективен.
